\documentclass[10pt]{article}
\usepackage{fullpage,amsfonts,mathpazo}
\usepackage{kotex}
\usepackage{mathtools}
\usepackage{tcolorbox}
\usepackage{caption}
\usepackage{subcaption}
\captionsetup[figure]{labelformat=empty}

\pagestyle{empty}

\begin{document}
\section*{ARA0111: Homework (Due: 12/21 23:59)}

소규모건축구조기준을 참고하여, 다음의 구조물 모델을 만들고 제작 보고서 제출

\begin{itemize}
	\item 제출물은 \texttt{팀명.zip} 파일이며, 아래의 파일들이 포함되어야 한다.   

\begin{enumerate}
	\item 최종보고서: \texttt{팀명.pdf}
	\item 사진: 3장 이상, \texttt{팀명\_01.png, 팀명\_02.png, ...}
\end{enumerate}
	\item 최종보고서의 본문은 표지 포함 10장을 넘기지 않도록 한다. 단, 워드나 한글파일로 제출될 경우 감점하며, 부록의 분량제한은 없다. 
	\item 최종보고서에는 일자별 회의장면이나 모델 제작과정이 포함되어야 한다. 
	\item 제작과정 회의에서는 구조적인 개념을 통한 제작 모델의 분석이 필요하며, 결론에 이러한 내용이 충분히 반영되어야 한다. 
	\item 구조물 선택은 임의로 하되 유명한 건물이 아니어도 상관없으며, 필요한 경우 가장 일반적인 소규모건축물구조기준에 제시된 상세가 활용된 1-Bay 건축물도 무관하다. 
	\item 평가기준은 세련된 구조물을 선택하고 이를 만들었냐는 전혀 중요하지 않다. 제작과정에 대한 고민이 충분히 이루어졌는가의 여부가 중요하며, 개선안이나 새로운 상세를 제시하면 가점이 부여된다. 
\end{itemize}
\end{document}
